% LaTeX/booktabs draft tables for the bounded slop-style paper package.
% Suggested preamble packages: booktabs, tabularx, threeparttable.
% These tables are submission-facing drafts; adjust widths/captions for the
% target venue template.

\begin{table*}[t]
\centering
\caption{Data and measurement scope.}
\label{tab:data-scope}
\begin{threeparttable}
\small
\begin{tabularx}{\textwidth}{@{}p{0.22\textwidth}X p{0.22\textwidth} p{0.20\textwidth}@{}}
\toprule
Measurement layer & Scope & Size & Paper role \\
\midrule
Dolma 3 pretraining reference & English-filtered retained non-code sample from a bounded scan & 5,740 documents; 8,689,472 tokens & Pretraining-style baseline \\
Dolci SFT targets & English-filtered retained target responses & 40,000 documents; 8,279,115 tokens & SFT data register \\
Dolci DPO pairs & English-filtered retained chosen/rejected pairs & 40,000 pairs; 80,000 response rows; 27,448,016 tokens & Preference-data contrast \\
Phase 1 pybiber & Full pybiber over retained Phase 1 samples & 125,740 documents $\times$ 67 features & Corpus-side register layer \\
OLMo reduced SGLang panels & Bounded free-running generation scope & 16,384 generations & OLMo realized-output and compounding layer \\
SmolLM3 no-think panel & Production-shape no-think generation grid & 49,152 generated records & Cross-ladder replication pressure \\
Phase 4 detector attribution & ModernBERT detector attribution pass & 10,000 attributed documents; 77,013 document-span rows & Candidate Tier-3 discovery \\
Phase 4 Tier-3 exact rerun & OLMo exact sequence-mass teacher-forced grid & 512 reference documents $\times$ 4 stages & Detector-candidate propensity check \\
\bottomrule
\end{tabularx}
\begin{tablenotes}
\footnotesize
\item Dolma counts refer to the retained English-filtered sample, not the full Dolma mixture. SmolLM3 generations are no-think mode for comparability.
\end{tablenotes}
\end{threeparttable}
\end{table*}

\begin{table*}[t]
\centering
\caption{Selected full-pybiber register means.}
\label{tab:pybiber-means}
\begin{threeparttable}
\small
\begin{tabularx}{\textwidth}{@{}p{0.20\textwidth}rrrrX@{}}
\toprule
Feature & Dolma & SFT & DPO chosen & DPO rejected & Interpretation \\
\midrule
Past tense & 57.065 & 4.315 & 6.421 & 8.805 & SFT sharply suppresses narrative time orientation \\
First-person pronouns & 44.732 & 9.360 & 9.223 & 11.934 & SFT/DPO are much less personal than Dolma \\
Third-person pronouns & 51.809 & 4.154 & 6.138 & 8.671 & SFT removes much narrative/person reference \\
Nominalizations & 8.944 & 24.805 & 27.373 & 28.913 & Post-training data are more nominalized \\
Attributive adjectives & 40.742 & 52.458 & 63.965 & 56.861 & Chosen responses are especially descriptive \\
Prepositions & 88.294 & 96.815 & 94.630 & 91.819 & SFT shifts toward expository relation marking \\
Adverbs & 73.226 & 31.683 & 42.427 & 37.954 & Dolma is more adverbial; DPO partially rebounds \\
Hedges & 0.742 & 0.057 & 0.163 & 0.174 & Broad shift is not simply more hedging \\
Amplifiers & 2.292 & 0.416 & 0.818 & 1.454 & Dolma exceeds SFT on intensification \\
Emphatics & 8.843 & 1.490 & 3.019 & 3.030 & Dolma exceeds SFT/DPO on emphatic marking \\
Clausal coordination & 12.089 & 3.004 & 4.912 & 5.189 & SFT is less coordination-heavy \\
\bottomrule
\end{tabularx}
\begin{tablenotes}
\footnotesize
\item Selected-feature document-bootstrap intervals are reported in Appendix Table~\ref{tab:pybiber-intervals}. DPO chosen/rejected contrasts are descriptive contrasts over retained Dolci preference pairs, not pure human-preference effects.
\end{tablenotes}
\end{threeparttable}
\end{table*}

\begin{table*}[t]
\centering
\caption{Tier-1 precision-validation status.}
\label{tab:tier1-precision}
\begin{threeparttable}
\small
\begin{tabularx}{\textwidth}{@{}p{0.21\textwidth}p{0.10\textwidth}p{0.12\textwidth}rrrrp{0.13\textwidth}X@{}}
\toprule
Feature & Role & Status & Labeled & Exact & False positive & Ambiguous & Precision & Paper treatment \\
\midrule
\texttt{contrastive\_negation} & Core & Pass & 61 & 57 & 3 & 0 & 0.934 & Interim precision-supported \\
\texttt{stock\_openers} & Core & Pass & 68 & 67 & 1 & 0 & 0.985 & Interim precision-supported \\
\texttt{stock\_closers} & Core & Pass & 50 & 48 & 2 & 0 & 0.960 & Interim precision-supported \\
\texttt{slop\_lexicon} & Core & Pass & 50 & 50 & 0 & 0 & 1.000 & Interim span/item precision-supported \\
\texttt{rule\_of\_three\_approx} & Core & Demote & 50 & 28 & 19 & 2 & 0.560 & Exploratory coordinated-triple surface only \\
\texttt{stock\_openers\_closers} & Derived & Derived & 0 & 0 & 0 & 0 & undefined & Treat as pooled convenience metric \\
\texttt{eqbench\_slop\_score} & Aggregate bridge & Unlabeled & 0 & 0 & 0 & 0 & undefined & Exploratory benchmark bridge \\
\texttt{slop\_lexicon\_v2\_candidate} & Candidate lexicon & Unlabeled & 0 & 0 & 0 & 0 & undefined & Exploratory candidate lexicon \\
\bottomrule
\end{tabularx}
\begin{tablenotes}
\footnotesize
\item The interim precision gate is 0.90 with enough labeled examples to support publication-grade regex claims. Passing precision with very few labels is not sufficient for headline use.
\end{tablenotes}
\end{threeparttable}
\end{table*}

\begin{table}[t]
\centering
\caption{Claim strength bands.}
\label{tab:claim-bands}
\small
\begin{tabularx}{\linewidth}{@{}p{0.34\linewidth}p{0.28\linewidth}X@{}}
\toprule
Band & Claim IDs & Allowed use \\
\midrule
Headline-ready, bounded & C1, C2, C7, C8, C9 & Main text with bounded wording \\
Paper-usable with explicit caveat & C3, C4, C5, C6, C10, C12, C13 & Main text only when caveat is colocated \\
Methods/status rather than substantive result & C11, C14 & Methods, limitations, or appendix \\
\bottomrule
\end{tabularx}
\end{table}

\begin{table*}[t]
\centering
\caption{Selected pybiber bootstrap intervals.}
\label{tab:pybiber-intervals}
\small
\begin{tabularx}{\textwidth}{@{}p{0.20\textwidth}rrrr@{}}
\toprule
Feature & Dolma & SFT & DPO chosen & DPO rejected \\
\midrule
Past tense & 57.065 [55.212, 59.078] & 4.315 [4.195, 4.460] & 6.421 [6.211, 6.618] & 8.805 [7.685, 9.963] \\
First-person pronouns & 44.732 [42.971, 46.451] & 9.360 [9.086, 9.792] & 9.223 [9.007, 9.449] & 11.934 [10.534, 13.503] \\
Third-person pronouns & 51.809 [50.260, 53.734] & 4.154 [4.022, 4.297] & 6.138 [5.950, 6.338] & 8.671 [7.820, 9.649] \\
Nominalizations & 8.944 [8.599, 9.303] & 24.805 [24.515, 25.119] & 27.373 [27.090, 27.642] & 28.913 [27.439, 31.086] \\
Attributive adjectives & 40.742 [40.094, 41.364] & 52.458 [51.821, 53.321] & 63.965 [63.474, 64.386] & 56.861 [55.840, 58.157] \\
Hedges & 0.742 [0.710, 0.772] & 0.057 [0.051, 0.064] & 0.163 [0.149, 0.177] & 0.174 [0.094, 0.265] \\
Amplifiers & 2.292 [2.226, 2.354] & 0.416 [0.400, 0.434] & 0.818 [0.795, 0.844] & 1.454 [0.550, 2.991] \\
Emphatics & 8.843 [8.696, 9.007] & 1.490 [1.447, 1.529] & 3.019 [2.971, 3.072] & 3.030 [2.760, 3.394] \\
\bottomrule
\end{tabularx}
\end{table*}

\begin{table*}[t]
\centering
\caption{Paper-safe negative claims.}
\label{tab:negative-claims}
\small
\begin{tabularx}{\textwidth}{@{}p{0.36\textwidth}X@{}}
\toprule
Avoided claim & Replacement \\
\midrule
DPO universally creates slop & Some markers amplify at specific stages and differ by ladder \\
Preference data are uniformly more slop-like & Dolci DPO chosen/rejected differences are feature-specific and register-specific \\
Alignment simply adds hedging & Pybiber shows a broader shift toward nominalized answer-register prose \\
EQ-Bench Slop Score is the outcome variable & EQ-Bench is an aggregate bridge; per-feature propensity remains the causal measurement layer \\
Detector clusters are human-perceived slop & Detector clusters are candidate Tier-3 features until human perceptibility labels exist \\
Reduced SGLang panels equal the original exhaustive grid & The paper estimand is the bounded production scope \\
\bottomrule
\end{tabularx}
\end{table*}
